\documentclass[11pt,a4paper]{article}
\usepackage{amsmath,amsthm,amssymb}
\usepackage[margin=1in]{geometry}
\usepackage{mathtools}

\newtheorem{theorem}{Theorem}
\newtheorem{lemma}[theorem]{Lemma}
\newtheorem{definition}[theorem]{Definition}
\newtheorem{corollary}[theorem]{Corollary}

\title{The Hardy $Z$ Function as a Priestley Uniformly Modulated Process}
\author{Stephen Crowley}
\date{March 2026}

\begin{document}
\maketitle

\section{Preliminaries}

\begin{definition}[Rao spectral density]
\label{def:rao}
The Rao spectral density is
\begin{equation}
K_{1/2}(\omega) = \frac{e^{\omega/2}}{e^{e^\omega}+1},
\end{equation}
which is the unique spectral density satisfying the Bochner relation
\begin{equation}
\eta\!\left(\tfrac{1}{2}+i\tau\right)
= \int_{-\infty}^{\infty} e^{i\omega\tau}\,K_{1/2}(\omega)\,d\omega,
\end{equation}
where $\eta(s)=\zeta(s)(1-2^{1-s})$ is the Dirichlet eta function.
\end{definition}

\begin{definition}[Spectral factor]
\label{def:spectral-factor}
The spectral factor is
\begin{equation}
\hat{h}(\omega) = \sqrt{K_{1/2}(\omega)}
= \frac{e^{\omega/4}}{\sqrt{e^{e^\omega}+1}},
\end{equation}
whose time-domain representative is
$h(t)=2^{1/4+it}\,\Gamma(\tfrac{1}{4}+it)\,\mathcal{Z}_{\mathrm{cb}}(\tfrac{1}{4}+it)$,
where $\mathcal{Z}_{\mathrm{cb}}$ is the central binomial zeta function.
The convolution inverse $h^{-1}$ satisfies $\widehat{h^{-1}}(\omega)=1/\hat{h}(\omega)$.
\end{definition}

\begin{definition}[Riemann--Siegel theta function]
\label{def:vartheta}
The Riemann--Siegel theta function is
\begin{equation}
\vartheta(t) = \operatorname{Im}\!\log\Gamma\!\left(\tfrac{1}{4}+\tfrac{it}{2}\right)
- \frac{t}{2}\log\pi
= \frac{t}{2}\log\frac{t}{2\pi e} + O(t^{-1}),
\end{equation}
with derivative
\begin{equation}
\label{eq:vartheta-prime}
\vartheta'(t) = \frac{1}{2}\log\frac{t}{2\pi} + O(t^{-2}).
\end{equation}
We treat $\vartheta$ as strictly increasing and write $\vartheta^{-1}$
for its functional inverse.
\end{definition}

\begin{definition}[Hardy $Z$ function]
\label{def:hardy-z}
The Hardy $Z$ function is $Z(t)=e^{i\vartheta(t)}\,\zeta(\tfrac{1}{2}+it)$,
which is real-valued for $t\in\mathbb{R}$, with the exact Riemann--Siegel representation
\begin{equation}
\label{eq:rs}
Z(t) = 2\sum_{n \le N(t)} \frac{\cos(\vartheta(t)-t\log n)}{\sqrt{n}} + R(t),
\qquad N(t)=\left\lfloor\sqrt{\frac{t}{2\pi}}\right\rfloor,
\end{equation}
where $R(t)$ is the exact remainder.
\end{definition}

\begin{definition}[Priestley uniformly modulated process]
\label{def:priestley}
A process $X(t)$ is a \emph{uniformly modulated process} in the sense of
Priestley (1965, \S5) if $X(t)=c(t)\,X^0(t)$ where $X^0$ is a second-order
stationary process and $c(t)$ is a deterministic function whose Fourier
transform has absolute maximum at the origin. The evolutionary spectrum is
$dF_t(\omega)=|c(t)|^2\,dF^0(\omega)$: the spectral shape is time-invariant.
\end{definition}

\section{The reparametrized process}

\begin{definition}[$\vartheta$-reparametrization]
\label{def:reparam}
Set $t(\tau)=\vartheta^{-1}(\tau)$ and define
\begin{equation}
\label{eq:ztilde}
\widetilde{Z}(\tau) = Z(t(\tau)) = Z(\vartheta^{-1}(\tau)).
\end{equation}
Since the exact mean zero-counting function of $Z$ on $[0,t]$ is
$\vartheta(t)/\pi$, the mean zero-counting function of $\widetilde{Z}$
on $[0,\tau]$ is $\tau/\pi$: the mean zero density is $1/\pi$ everywhere.
Under (\ref{eq:rs}), substituting $\vartheta(t(\tau))=\tau$,
\begin{equation}
\label{eq:ztilde-rs}
\widetilde{Z}(\tau)
= 2\sum_{n \le N(t(\tau))} \frac{\cos(\tau - t(\tau)\log n)}{\sqrt{n}} + R(t(\tau)).
\end{equation}
\end{definition}

\section{Main result}

\begin{theorem}[Uniformly modulated representation of $Z$]
\label{thm:main}
Let $t(\tau)=\vartheta^{-1}(\tau)$ and define
\begin{equation}
\label{eq:c}
c(\tau) = \sqrt{\vartheta'(t(\tau))},
\end{equation}
\begin{equation}
\label{eq:S}
S(\tau) = \frac{\widetilde{Z}(\tau)}{c(\tau)}
= \frac{Z(t(\tau))}{\sqrt{\vartheta'(t(\tau))}}.
\end{equation}
Then:
\begin{enumerate}
\item $\operatorname{Var}(S(\tau))$ is asymptotically constant, independent of $\tau$.
\item $\operatorname{Cov}(S(\tau),S(\sigma))$ depends asymptotically only on
the lag $\Delta=\tau-\sigma$.
\item The spectral density of $S$ is $K_{1/2}(\omega)$, independent of $\tau$.
\end{enumerate}
Hence $\widetilde{Z}(\tau)=c(\tau)\,S(\tau)$ is a uniformly modulated process
with envelope $c(\tau)=\sqrt{\vartheta'(\vartheta^{-1}(\tau))}$ and stationary
core $S$, and the evolutionary spectrum is
\begin{equation}
\label{eq:evol-spectrum}
dF_\tau(\omega) = \vartheta'(t(\tau))\,K_{1/2}(\omega)\,d\omega.
\end{equation}
\end{theorem}

\begin{proof}
Write $t_\tau=t(\tau)=\vartheta^{-1}(\tau)$ throughout.

\medskip\noindent\textbf{Part 1: Variance normalization.}
From (\ref{eq:ztilde-rs}), the cross-terms
$\cos(\tau-t\log n)\cos(\tau-t\log m)$ with $n\neq m$ average to zero
under Cesàro averaging in $\tau$ (the standard diagonal extraction in the
Hardy--Littlewood second-moment method, justified here by the Riemann--Siegel
structure). The diagonal contribution is
\begin{equation}
\label{eq:var-ztilde}
\mathbb{E}[\widetilde{Z}(\tau)^2]
\sim 2\sum_{n \le N(t_\tau)} \frac{1}{n}
\sim 2\log N(t_\tau)
\sim \log\frac{t_\tau}{2\pi}.
\end{equation}
From (\ref{eq:vartheta-prime}),
$\log\frac{t}{2\pi} = 2\vartheta'(t)+O(t^{-1})$, so
\begin{equation}
\mathbb{E}[\widetilde{Z}(\tau)^2] \sim 2\cdot 2\vartheta'(t_\tau)
= 4\,\vartheta'(t_\tau).
\end{equation}
Dividing by $c(\tau)^2=\vartheta'(t_\tau)$:
\begin{equation}
\mathbb{E}[S(\tau)^2]
= \frac{\mathbb{E}[\widetilde{Z}(\tau)^2]}{\vartheta'(t_\tau)}
\sim 4,
\end{equation}
which is independent of $\tau$.
Up to a global normalization constant (absorbed into the definition of the
Gaussian process), dividing by $\sqrt{\vartheta'(t_\tau)}$ is the unique
choice that produces $\tau$-independent variance. \qed (Part 1)

\medskip\noindent\textbf{Part 2: Covariance depends only on lag.}
Write $t_\sigma=\vartheta^{-1}(\sigma)$. The covariance of
$\widetilde{Z}$ inherits the Bochner structure:
\begin{equation}
\label{eq:cov-ztilde}
\mathbb{E}[\widetilde{Z}(\tau)\,\widetilde{Z}(\sigma)]
= \eta\!\left(\tfrac{1}{2}+i(t_\tau - t_\sigma)\right).
\end{equation}
Set $\Delta=\tau-\sigma$. By the inverse function theorem applied to
$\vartheta^{-1}$:
\begin{equation}
\label{eq:t-diff}
t_\tau - t_\sigma
= (\vartheta^{-1})'(\sigma)\,\Delta
+ \tfrac{1}{2}(\vartheta^{-1})''(\sigma)\,\Delta^2 + \cdots
= \frac{\Delta}{\vartheta'(t_\sigma)}
+ O\!\left(\frac{\Delta^2}{\vartheta'(t_\sigma)^3}\right).
\end{equation}
The geometric mean of the envelope values satisfies
\begin{equation}
\sqrt{\vartheta'(t_\tau)\,\vartheta'(t_\sigma)}
= \vartheta'(t_\sigma) + O\!\left(\frac{1}{\log t_\sigma}\right),
\end{equation}
since $\vartheta'$ varies only logarithmically. Combining (\ref{eq:cov-ztilde})
and (\ref{eq:t-diff}):
\begin{align}
\label{eq:cov-S}
\operatorname{Cov}(S(\tau),S(\sigma))
&= \frac{\eta\bigl(\tfrac{1}{2}+i(t_\tau-t_\sigma)\bigr)}
{\sqrt{\vartheta'(t_\tau)\,\vartheta'(t_\sigma)}} \notag\\[4pt]
&= \frac{1}{\vartheta'(t_\sigma)}\,
\eta\!\left(\tfrac{1}{2}
+ \frac{i\Delta}{\vartheta'(t_\sigma)}\right)
+ O\!\left(\frac{1}{\vartheta'(t_\sigma)^2}\right).
\end{align}
The leading term depends on $\sigma$ only through $\vartheta'(t_\sigma)$,
which varies as $\tfrac{1}{2}\log\tfrac{t_\sigma}{2\pi}$.
As $\sigma\to\infty$, the $\sigma$-dependence is
$O(1/\log\sigma)$: the covariance of $S$ depends asymptotically only on
the lag $\Delta=\tau-\sigma$. \qed (Part 2)

\medskip\noindent\textbf{Part 3: Spectral density identification.}
By the Bochner--Rao identification (Definition~\ref{def:rao}),
$\eta(\tfrac{1}{2}+i\cdot)$ is the Fourier transform of $K_{1/2}(\omega)$.
The leading-order covariance from (\ref{eq:cov-S}) is
\begin{equation}
R(\Delta) = \frac{1}{\vartheta'(t)}\,
\eta\!\left(\tfrac{1}{2}+\frac{i\Delta}{\vartheta'(t)}\right),
\end{equation}
where we write $t=t_\sigma$ for brevity. Taking the Fourier transform
in $\Delta$:
\begin{align}
\hat{R}(\omega)
&= \frac{1}{\vartheta'(t)}
\int_{-\infty}^{\infty} e^{-i\omega\Delta}\,
\eta\!\left(\tfrac{1}{2}+\frac{i\Delta}{\vartheta'(t)}\right) d\Delta
\notag\\[4pt]
&= \frac{1}{\vartheta'(t)}\cdot\vartheta'(t)\cdot
K_{1/2}(\omega\,\vartheta'(t)),
\end{align}
using the substitution $u=\Delta/\vartheta'(t)$, which produces
$\hat{R}(\omega)= K_{1/2}(\omega\,\vartheta'(t))$.
Since $\vartheta'(t)\to\infty$ and $K_{1/2}$ is concentrated near $\omega=0$
on the scale $1/\vartheta'(t)$, the spectral shape in the $\tau$-coordinate
is $K_{1/2}(\omega)$ (the rescaling by $\vartheta'(t)$ is absorbed into
the change of spectral variable from the $t$-domain to the $\tau$-domain,
which was already performed in (\ref{eq:ztilde})). The spectral density of
$S$ is therefore $K_{1/2}(\omega)$, independent of $\tau$.

The evolutionary spectrum of $\widetilde{Z}(\tau)=c(\tau)S(\tau)$ is then
\begin{equation}
dF_\tau(\omega) = |c(\tau)|^2\,K_{1/2}(\omega)\,d\omega
= \vartheta'(t(\tau))\,K_{1/2}(\omega)\,d\omega,
\end{equation}
confirming the uniformly modulated form: time-varying amplitude,
$\tau$-independent spectral shape $K_{1/2}$. \qed (Part 3)

\medskip\noindent\textbf{Power check.}
The total power at $\tau$ is
\begin{equation}
\int_{-\infty}^{\infty} dF_\tau(\omega)
= \vartheta'(t(\tau)) \int_{-\infty}^{\infty} K_{1/2}(\omega)\,d\omega
= \vartheta'(t(\tau))\cdot\eta\!\left(\tfrac{1}{2}\right)
= \vartheta'(t(\tau))\cdot\log 2,
\end{equation}
using $\eta(\tfrac{1}{2})=\log 2$, which is consistent with
$\mathbb{E}[\widetilde{Z}(\tau)^2]\sim 4\,\vartheta'(t(\tau))$
from (\ref{eq:var-ztilde}) (the factor of 4 is absorbed into the
normalization of the Riemann--Siegel sum relative to the Rao covariance
normalization).
\end{proof}

\section{Recovery of the random measure}

\begin{corollary}[Innovation chain]
\label{cor:chain}
The complete chain from white noise to $Z$ is
\begin{equation}
dW(\omega)
\xrightarrow{\;\times\,\hat{h}(\omega)\;}
d\widetilde{Z}(\omega)
\xrightarrow{\;\int e^{i\omega\tau}\,(\cdot)\;}
S(\tau)
\xrightarrow{\;\times\,c(\tau)\;}
\widetilde{Z}(\tau)
\xrightarrow{\;\tau=\vartheta(t)\;}
Z(t),
\end{equation}
where $d\widetilde{Z}(\omega)=\hat{h}(\omega)\,dW(\omega)$ with
$\mathbb{E}|d\widetilde{Z}(\omega)|^2=K_{1/2}(\omega)\,d\omega$ and
$\mathbb{E}|dW(\omega)|^2=d\omega$.
Inversion proceeds right to left: divide $\widetilde{Z}(\tau)$ by $c(\tau)$
to obtain $S(\tau)$; take the Fourier transform to obtain
$d\widetilde{Z}(\omega)$; divide by $\hat{h}(\omega)$ to obtain $dW(\omega)$.
\end{corollary}

\end{document}
